\thispagestyle{empty}

\vspace*{1cm}

\begin{center}
    \textbf{Abstract}
\end{center}

\vspace*{1cm}

\noindent 
Der zunehmende Einsatz künstlicher Intelligenz im Coaching eröffnet neue Möglichkeiten für die digitale Unterstützung persönlicher Entwicklungs- und Reflexionsprozesse. Insbesondere im traumasensiblen Coaching ergeben sich jedoch besondere Anforderungen an die Gestaltung von KI-basierten Dialogsystemen, da emotionale Sicherheit, Selbstregulation und eine situationsgerechte Gesprächsführung gewährleistet werden müssen. Vor diesem Hintergrund untersucht die vorliegende Arbeit, inwieweit eine theoriegeleitete, promptbasierte Dialogsteuerungslogik die Funktionalität eines KI-basierten Coaching-Chatbots verbessern und eine traumasensible Unterstützung ermöglichen kann.

Ziel der Arbeit ist die Entwicklung und Evaluation eines psychologisch fundierten Chatbot-Frameworks auf Grundlage der Persönlichkeits-System-Interaktionen-Theorie (PSI-Theorie) nach Kuhl. Im Rahmen eines Design-Science-Research-Ansatzes werden zentrale psychologische Wirkmechanismen in eine strukturierte technische Entscheidungslogik überführt. Das entwickelte Framework verbindet eine mehrstufige Sicherheitsprüfung mit einer PSI-basierten Zustands- und Potenzialanalyse sowie einer zielgerichteten Auswahl und Steuerung von Interventionen. Die Evaluation erfolgt anhand von 400 kontrollierten Dialogsimulationen, in denen die PSI-basierte Dialogsteuerung einer konventionellen, traumasensiblen Kontrollbedingung gegenübergestellt wird. Die statistische Auswertung basiert auf dem einseitigen Wilcoxon-Vorzeichen-Rang-Test.

Die Ergebnisse zeigen unterschiedliche Effekte hinsichtlich der beiden untersuchten Modulationsannahmen. Für die Willensbahnung konnte durch die PSI-basierte Steuerungslogik eine signifikant stärkere Aktivierung der intuitiven Verhaltenssteuerung gegenüber der Kontrollbedingung festgestellt werden (p < 0,001; r = 0,872). Für die Selbstentwicklung durch Affektregulation ließ sich hingegen keine statistisch signifikante Überlegenheit nachweisen. Die Ergebnisse verdeutlichen somit, dass eine theoriegeleitete Dialogsteuerung nicht für alle psychologischen Wirkmechanismen gleichermaßen wirksam ist.

Die Arbeit zeigt die grundsätzliche technische Umsetzbarkeit einer PSI-basierten Dialogsteuerung und leistet einen konzeptionellen Beitrag zur Gestaltung psychologisch fundierter KI-Coaching-Systeme. Eine tatsächliche psychologische oder therapeutische Wirksamkeit lässt sich aufgrund der ausschließlich simulationsbasierten Untersuchung jedoch nicht ableiten. Diese bedarf einer weiterführenden empirischen Überprüfung mit realen Nutzerinnen und Nutzern.

**Schlüsselwörter:** Künstliche Intelligenz, Coaching-Chatbots, traumasensibles Coaching, PSI-Theorie, Prompt Engineering, Dialogsteuerung, Selbstregulation
